\documentclass[11pt]{article}
\usepackage[spanish]{babel}
\usepackage[utf8]{inputenc}
\usepackage[T1]{fontenc}
\usepackage{geometry}
\geometry{margin=2.5cm}
\usepackage{amsmath}
\usepackage{enumitem}
\usepackage{booktabs}
\usepackage{array}
\usepackage{natbib}
\usepackage{hyperref}
\hypersetup{colorlinks=true, linkcolor=blue, citecolor=blue, urlcolor=blue}

\title{Estrategia de identificación\\
\large Shocks de precios de metales y estructura del empleo en el Perú:\\
un diseño \textit{shift-share} (Bartik) centrado en la factibilidad}
\author{}
\date{Julio 2026}

\begin{document}
\maketitle

\begin{abstract}
\noindent Este documento propone una estrategia de identificación \emph{factible} para
estimar el efecto causal de los shocks de precios internacionales de metales sobre la
estructura sectorial y de formalidad del empleo de los hogares en las regiones peruanas.
La estrategia es un diseño \textit{shift-share} (Bartik): la exposición regional se construye
combinando la especialización minera \emph{pre-determinada} de cada región con la variación
\emph{exógena} de los precios internacionales de cada metal. El énfasis está puesto en tres
cosas prácticas: (i) que todos los datos necesarios son públicos y descargables; (ii) que el
supuesto de identificación es creíble porque el Perú —y con más razón una sola de sus
regiones— es tomador de precios en los mercados mundiales de metales; y (iii) que los
supuestos son \emph{contrastables} con pruebas estándar. Se incluye una checklist de
factibilidad, las ecuaciones a estimar, las amenazas y sus remedios, y un plan mínimo
viable (MVP) para una primera versión del trabajo.
\end{abstract}

\tableofcontents

\section{Idea en una frase y por qué es identificable}
Queremos comparar regiones peruanas que, por su composición productiva histórica, están
\emph{más} o \emph{menos} expuestas a los metales cuyos precios internacionales suben o bajan.
La variación explotada no es ``tener minería'' (endógena a muchas cosas), sino la
\textbf{interacción} entre una exposición fija de base y un precio mundial que ninguna región
peruana puede mover. Esa interacción es la fuente de identificación.

\medskip
\noindent\textbf{El activo central de credibilidad:} los precios internacionales de cobre, oro
y zinc se determinan en mercados globales (LME, mercados de futuros). Una región del Perú
—Áncash, Cajamarca, Arequipa— es un actor despreciable en ellos. Por tanto, condicional a la
composición productiva de base, el shock de precios es plausiblemente \emph{exógeno} al
mercado laboral de esa región. Este es exactamente el supuesto de ``exogeneidad de los
\textit{shifts}'' de \citet{borusyak2022}, y es lo que hace el diseño defendible.

\section{Datos: todo es público (factibilidad de acceso)}

La Tabla~\ref{tab:datos} resume las fuentes. El punto clave es que \emph{no se requiere ningún
dato confidencial ni convenios}: todo está disponible en portales oficiales.

\begin{table}[h]
\centering
\small
\caption{Fuentes de datos, todas públicas}
\label{tab:datos}
\begin{tabular}{p{3.2cm} p{4.2cm} p{2.3cm} p{4.4cm}}
\toprule
\textbf{Componente} & \textbf{Variable} & \textbf{Unidad/frec.} & \textbf{Fuente} \\
\midrule
Empleo, ingreso, informalidad & Sector, categoría ocupacional, formalidad (afiliación/registro), ingreso laboral & Individuo/hogar; anual & ENAHO (INEI), microdatos abiertos \\
\addlinespace
Precios de metales (\textit{shifts}) & Precio internacional de cobre, oro, zinc, plomo, plata & Mensual/anual & Banco Mundial ``Pink Sheet''; BCRP (series estadísticas); LME \\
\addlinespace
Exposición minera (\textit{shares}) & VAB minero o empleo minero por región en año base & Región/depto.; base fija & INEI (cuentas nacionales departamentales); Censo/ENAHO del año base \\
\addlinespace
Producción física por metal & Toneladas/onzas por unidad minera y distrito & Distrito; anual & MINEM (Anuario/boletines); INGEMMET (concesiones) \\
\addlinespace
Canal fiscal (control) & Transferencias de canon minero y gasto ejecutado & Distrito/región; anual & MEF -- Consulta Amigable (SIAF) \\
\addlinespace
Deflactor & IPC EE.UU.\ (precios reales) e IPC Perú & Mensual & BLS / INEI \\
\bottomrule
\end{tabular}
\end{table}

\paragraph{Nivel geográfico (decisión de factibilidad).} La ENAHO es representativa a nivel de
\textbf{departamento} (25 regiones) todos los años, y en varios años permite mayor
desagregación. Recomiendo fijar la unidad de exposición $r$ en el \textbf{departamento} para
garantizar representatividad y un panel balanceado de regiones a lo largo de 2004--2023. La
producción minera por distrito (MINEM) se agrega al departamento para construir las
participaciones. Usar el departamento evita el principal problema de factibilidad de bajar a
provincia/distrito: la pérdida de representatividad muestral de la ENAHO.

\paragraph{Estructura temporal.} Se puede usar la ENAHO como \textbf{corte transversal
repetido} (2004--2023), que es lo más robusto en tamaño muestral. El componente de
\textbf{panel rotativo} de la ENAHO (hogares seguidos, disponible desde $\sim$2007 en paneles
de 2--5 años) puede usarse como extensión para mirar transiciones individuales
formal$\leftrightarrow$informal, pero \emph{no} debe ser la base principal por su menor tamaño
y atrición. Recomendación: identificar con el corte repetido; usar el panel como chequeo de
mecanismo.

\section{Construcción del instrumento \textit{shift-share}}

Para cada región $r$ y año $t$ definimos el índice de precios relevante para esa región:
\begin{equation}
B_{rt} \;=\; \sum_{m \in M} s_{rm}^{\,0}\; \Delta \ln p_{mt},
\label{eq:bartik}
\end{equation}
donde:
\begin{itemize}[leftmargin=1.4em, itemsep=1pt]
  \item $M=\{$cobre, oro, zinc, plomo, plata$\}$ (los metales de mayor peso exportador);
  \item $s_{rm}^{\,0}$ es la \textbf{participación pre-determinada} del metal $m$ en la
  economía minera de la región $r$, medida en un \textbf{año base} $0$ anterior a la ventana
  de análisis (p.\ ej.\ producción física del metal $m$ en $r$ sobre producción minera total
  de $r$, o su participación en el VAB/empleo);
  \item $\Delta \ln p_{mt}$ es el cambio (log) del \textbf{precio internacional real} del
  metal $m$, común a todas las regiones.
\end{itemize}

\paragraph{Por qué esta construcción es la correcta (y factible).} Las participaciones
$s_{rm}^{0}$ se fijan \emph{antes} del periodo de estimación, de modo que la variación
año-a-año de $B_{rt}$ proviene solo de los precios mundiales. Esto materializa la separación
entre ``dónde estabas expuesto'' (share, potencialmente endógeno pero fijo) y ``qué pasó con
el precio'' (shift, exógeno). Con 5 metales y $\sim$20 departamentos mineros se obtiene amplia
variación transversal en la exposición.

\paragraph{Variante recomendada para robustez.} Construir también $B_{rt}$ con participaciones
del \emph{empleo} minero (no solo producción física), y una versión con precios
\emph{rezagados} un periodo, para separar efectos contemporáneos de dinámicos.

\section{Especificación econométrica}

\subsection{Ecuación principal (niveles del resultado)}
Para el individuo $i$ en la región $r$ y año $t$:
\begin{equation}
y_{irt} \;=\; \beta\, B_{rt} \;+\; \mathbf{X}_{irt}'\gamma \;+\; \alpha_r \;+\; \delta_t
\;+\; \big(\alpha_r \cdot t\big) \;+\; \varepsilon_{irt},
\label{eq:main}
\end{equation}
con efectos fijos de región $\alpha_r$ (absorben todo lo permanente de cada departamento:
geografía, dotación, instituciones), efectos fijos de año $\delta_t$ (absorben el ciclo
macro-nacional y los shocks comunes), tendencias lineales por región $\alpha_r\cdot t$
(controlan trayectorias divergentes preexistentes), y controles individuales $\mathbf{X}$
(edad, sexo, educación, urbano/rural). El coeficiente de interés es $\beta$.

\paragraph{Resultados $y_{irt}$.}
\begin{enumerate}[label=(\alph*), leftmargin=1.6em, itemsep=1pt]
  \item \textbf{Informalidad}: indicador $=1$ si el empleo es informal (LPM), interpretación
  directa de $\beta$ como cambio en la probabilidad.
  \item \textbf{Estructura sectorial}: modelo \textbf{multinomial} sobre las celdas
  (minería / otro transable / no transable) $\times$ (formal / informal), para ver la
  reasignación completa.
  \item \textbf{Ingreso laboral} (log) como resultado secundario/mecanismo.
\end{enumerate}

\subsection{Asimetría auge vs.\ caída (contribución central)}
Separamos la parte positiva y negativa del shock:
\begin{equation}
y_{irt} \;=\; \beta^{+} B_{rt}^{+} \;+\; \beta^{-} B_{rt}^{-} \;+\;
\mathbf{X}_{irt}'\gamma + \alpha_r + \delta_t + (\alpha_r\cdot t) + \varepsilon_{irt},
\end{equation}
donde $B_{rt}^{+}=\max(B_{rt},0)$ y $B_{rt}^{-}=\min(B_{rt},0)$. El test $\beta^{+}\neq
|\beta^{-}|$ evalúa si la informalidad amortigua las caídas más de lo que se reduce en los
auges. \emph{Esto es lo que la literatura peruana no ha estimado y es barato de implementar.}

\subsection{Aislar el canal laboral del canal fiscal (canon)}
Para que $\beta$ capture mercado laboral y no gasto público inducido por canon, se incluye la
transferencia de canon per cápita $C_{rt}$ como control (o como mediador en una descomposición):
\begin{equation}
y_{irt} \;=\; \beta\, B_{rt} \;+\; \theta\, C_{rt} \;+\; \mathbf{X}_{irt}'\gamma
+ \alpha_r + \delta_t + (\alpha_r\cdot t) + \varepsilon_{irt}.
\end{equation}
Comparar $\beta$ con y sin $C_{rt}$ indica cuánto del efecto opera por la vía fiscal. Los
datos de canon están en la Consulta Amigable del MEF (factible).

\section{Supuesto de identificación y cómo se contrasta}

\paragraph{Supuesto.} Condicional a $\alpha_r,\delta_t$ y las tendencias, la variación de
$B_{rt}$ es tan buena como aleatoria respecto a $\varepsilon_{irt}$. Bajo el marco de
\citet{borusyak2022} basta con la \textbf{exogeneidad de los \textit{shifts}} (precios
mundiales), que aquí es fuerte: son precios globales frente a una economía regional
despreciable. Complementariamente, \citet{goldsmith2020} justifican el diseño desde la
exogeneidad de las \emph{participaciones} de base.

\paragraph{Contrastes de credibilidad (todos estándar y factibles).}
\begin{enumerate}[leftmargin=1.6em, itemsep=1pt]
  \item \textbf{Pre-tendencias}: regresión de \emph{leads} del shock; si $B_{r,t+1}$ ``predice''
  resultados en $t$, hay problema. Fácil de graficar (event-study).
  \item \textbf{Balance de las participaciones}: correlacionar $s_{rm}^{0}$ con covariables
  regionales pre-base (\citealp{goldsmith2020}); reportar los pesos de Rotemberg para saber
  qué metal/región manda la estimación.
  \item \textbf{Inferencia correcta}: errores estándar agrupados por región y por año; y/o
  inferencia a nivel de \emph{shock} (AKM) de \citet{adao2019}, que corrige la correlación
  entre regiones con exposición parecida.
  \item \textbf{Placebos}: usar precios de metales que la región no produce como falso shock.
  \item \textbf{Robustez de base}: variar el año base y la definición de participación
  (producción vs.\ empleo).
\end{enumerate}

\section{Amenazas a la identificación y remedios (factibles)}

\begin{table}[h]
\centering
\small
\caption{Amenazas y cómo se manejan sin datos adicionales}
\begin{tabular}{p{4.6cm} p{9.6cm}}
\toprule
\textbf{Amenaza} & \textbf{Remedio} \\
\midrule
Migración endógena (los trabajadores se mueven hacia regiones en auge, sesgando la muestra) &
Usar región de \emph{residencia habitual}/nacimiento cuando esté disponible; reportar efectos
sobre composición demográfica como chequeo; interpretar como efecto sobre residentes. \\
\addlinespace
Participaciones endógenas (la especialización minera no es aleatoria) & Fijarlas en año base
pre-muestra; el FE de región absorbe su nivel; identificación viene del \textit{shift}
(\citealp{borusyak2022}). \\
\addlinespace
Precios de metales correlacionados con el ciclo mundial (demanda de China) & Los FE de año
$\delta_t$ absorben todo shock \emph{común}; $\beta$ usa solo variación \emph{diferencial} por
exposición. \\
\addlinespace
Pocos ``clusters'' (25 departamentos) para inferencia & Inferencia a nivel de shock (AKM,
\citealp{adao2019}); \textit{wild cluster bootstrap} como respaldo. \\
\addlinespace
Confusión con canon/gasto público & Controlar por canon per cápita del MEF; reportar $\beta$
con y sin control. \\
\bottomrule
\end{tabular}
\end{table}

\section{Checklist de factibilidad}

\begin{enumerate}[label=$\square$, leftmargin=1.6em, itemsep=2pt]
  \item Microdatos ENAHO 2004--2023 descargables del INEI (formato abierto). \textbf{Sí.}
  \item Precios de metales (Banco Mundial Pink Sheet / BCRP). Series largas, gratuitas. \textbf{Sí.}
  \item Producción minera por distrito/metal (MINEM). Anuarios públicos. \textbf{Sí.}
  \item Transferencias de canon (MEF Consulta Amigable). Público. \textbf{Sí.}
  \item Software: cualquier stack (R/Stata/Python); paquetes de shift-share e inferencia AKM
  disponibles y gratuitos. \textbf{Sí.}
  \item Sin datos confidenciales, sin convenios, sin trabajo de campo. \textbf{Sí.}
\end{enumerate}

\noindent\textbf{Cuello de botella real} (no de acceso, sino de trabajo): la
\emph{armonización} de la clasificación sectorial y de la definición de informalidad de la
ENAHO a lo largo de 20 años (cambios de CIIU y de cuestionario). Es tedioso pero estándar;
existe documentación del INEI y literatura previa que lo ha hecho.

\section{Plan mínimo viable (MVP) para la primera versión}
\begin{enumerate}[leftmargin=1.6em, itemsep=1pt]
  \item Construir $s_{rm}^{0}$ con producción física por departamento en el año base y
  $\Delta\ln p_{mt}$ del Pink Sheet; formar $B_{rt}$.
  \item Estimar la ecuación~\eqref{eq:main} con \textbf{informalidad} como único resultado y
  FE de región y año. Reportar el event-study de pre-tendencias.
  \item Añadir la descomposición auge/caída y el control de canon.
  \item Extender al modelo multinomial de estructura del empleo.
\end{enumerate}
Los pasos 1--2 ya constituyen un resultado publicable como \emph{stylized fact} causal; los
pasos 3--4 son la contribución diferencial.

\section{Conclusión sobre factibilidad}
La estrategia es \textbf{altamente factible}: (i) datos 100\% públicos y gratuitos; (ii) un
supuesto de identificación fuerte y bien fundamentado (precios mundiales exógenos a una región
peruana); (iii) supuestos \emph{contrastables} con herramientas estándar; y (iv) un camino
incremental (MVP $\to$ contribución) que permite tener un primer resultado creíble con
relativamente poco. El riesgo principal no es de identificación ni de acceso a datos, sino de
\emph{limpieza y armonización} de la ENAHO, que es trabajo conocido y acotado.

\begin{thebibliography}{9}
\setlength{\itemsep}{2pt}
\bibitem[Adão et al.(2019)]{adao2019}
Adão, R., Kolesár, M. y Morales, E. (2019). Shift-Share Designs: Theory and Inference.
\textit{Quarterly Journal of Economics}, 134(4), 1949--2010.
\bibitem[Borusyak et al.(2022)]{borusyak2022}
Borusyak, K., Hull, P. y Jaravel, X. (2022). Quasi-Experimental Shift-Share Research Designs.
\textit{Review of Economic Studies}, 89(1), 181--213.
\bibitem[Goldsmith-Pinkham et al.(2020)]{goldsmith2020}
Goldsmith-Pinkham, P., Sorkin, I. y Swift, H. (2020). Bartik Instruments: What, When, Why, and
How. \textit{American Economic Review}, 110(8), 2586--2624.
\end{thebibliography}

\end{document}
